\section{Introduction}
Tuning the position PID of a geared DC-motor joint is usually done by hand on hardware, starting from a manufacturer datasheet that is incomplete and sometimes internally inconsistent. Reinforcement learning (RL) has been proposed to tune or adapt PID gains automatically \cite{carlucho2017incremental,lawrence2021shallow}, but evaluations are often limited to simulation, to few random seeds and to weak baselines \cite{henderson2018matters,agarwal2021precipice}. A fair question is therefore not whether a learned tuner beats a poorly chosen preset, but how much of its apparent gain survives a properly tuned \emph{constant} baseline and a simple non-learned scenario-conditioned rule, and how much of what any scenario-aware tuner could achieve it captures.

\paragraph*{What this paper does.}
On an open digital twin of a geared brushed DC-motor joint with a payload arm (Section~\ref{sec:problem}) we compare fixed and one-shot tuners with PPO, SAC and TD3 policies that adapt the gains every \SI{20}{ms} (Section~\ref{sec:method}). Every comparison is paired on the same scenarios, baselines and the reported policy seed are chosen on separate validation seeds, and the analysis rules (primary metric, one primary test, correction for the others, a separate table for infeasible scenarios) are stated before the confirmatory seeds are touched (Section~\ref{sec:setup}).

\paragraph*{What it finds.}
(i) An unconstrained multiplicative PPO setup failed against the validation-tuned constant and its training was unstable (Section~\ref{sec:v2}). (ii) A residual formulation, in which the action is an offset from that constant, improved on it modestly and consistently across training seeds and on fresh scenarios (Section~\ref{sec:v3}), but captured only part of the oracle headroom. (iii) SAC and TD3 did not reproduce the benefit (Section~\ref{sec:v4}). (iv) The learned policies end at gains that vary little across scenarios; a post-hoc check on further fresh seeds shows that a much larger validation search for the constant recovers most of the gain and that a simple non-learned load-conditioned rule matches the policy (Section~\ref{sec:posthoc}). The headline is therefore the protocol, and the residual-RL result is secondary.

\paragraph*{Scope.}
Everything is simulated on a datasheet-derived twin that has not been compared with hardware; one motor model and one reward design are used, with three seeds per algorithm. The claims are therefore pre-validation, simulation-only claims. The hardware experiments that would test them are described as a replication protocol in Appendix~\ref{app:hardware} and contain no results.

\paragraph*{Contributions} (only those the evidence supports): a reproducible evaluation protocol for learned PID adaptation (validation-selected baselines with a generous constant search, a non-learned scenario-conditioned baseline, exploratory and confirmatory seed sets, an oracle bound, infeasible scenarios reported separately); the evidence that on this twin most of an apparent residual-PPO gain over a modestly tuned constant is available without learning; and negative results (an unconstrained setup, TD3, a mixed SAC) with the code, data and policies to check them.
